\chapter{Hyperventilation}

\section*{Definition}
Hyperventilation is rapid or deep breathing that exceeds the body's metabolic demand, leading to excessive elimination of carbon dioxide (hypocapnia) and respiratory alkalosis. It is most commonly triggered by anxiety but can occur with various medical conditions.

\section*{What Happens in Your Body}
Ventilation removes CO2 from the blood. When ventilation exceeds metabolic CO2 production, arterial CO2 drops and pH rises (respiratory alkalosis). Alkalosis causes cerebral narrowing of blood vessels, reduced cerebral blood flow, and altered neuronal excitability, producing lightheadedness, paresthesias, and muscle cramps. Alkalosis also increases binding of calcium to albumin, reducing ionized calcium and causing neuromuscular hyperexcitability (Chvostek and Trousseau signs).

\section*{Causes}
\begin{itemize}
\item Anxiety, panic attacks, or acute stress (most common cause)
\item Pulmonary causes: asthma exacerbation, pulmonary embolism, pneumonia
\item Cardiac causes: heart failure, myocardial infarction
\item Metabolic acidosis compensation (diabetic ketoacidosis)
\item Salicylate overdose (aspirin toxicity)
\item Pregnancy (progesterone stimulates respiration)
\item Liver cirrhosis
\item High altitude
\item Central nervous system disorders
\item Fever or sepsis
\end{itemize}

\section*{When to See a Doctor}
\begin{itemize}
\item Rapid, deep breathing out of proportion to exertion
\item Lightheadedness, dizziness, or feeling faint
\item Tingling or numbness of the lips, fingers, and toes
\item Carpopedal spasm: muscle cramping in hands and feet
\item Chest tightness or palpitations
\item Respiratory alkalosis on arterial blood gas
\item Anxiety or impending sense of doom
\item Difficulty speaking in full sentences
\end{itemize}

\section*{Related Symptoms}
\begin{itemize}
\item Anxiety and panic attacks
\item Shortness of breath (dyspnea)
\item Chest pain
\item Palpitations
\item Dizziness and presyncope
\item Tingling and numbness (paresthesias)
\item Hypocalcemia symptoms (tetany)
\end{itemize}

\section*{Self-Care / Home Management}
\begin{itemize}
\item Recognize hyperventilation and practice slow, controlled breathing
\item Inhale deeply through the nose for 4 seconds, hold for 4 seconds, exhale slowly through pursed lips for 6 seconds
\item Reassure yourself that the episode is temporary and not dangerous
\item Do NOT breathe into a paper bag (can cause hypoxia)
\item Remove yourself from stressful situations if possible
\item Practice relaxation techniques and mindfulness
\end{itemize}

\section*{How Your Doctor Will Evaluate You}
\begin{itemize}
\item Clinical history: acute onset with anxiety or stress
\item Observation of breathing pattern
\item Arterial blood gas shows respiratory alkalosis and low PaCO2
\item Electrolytes, calcium, magnesium
\item Chest X-ray if pulmonary cause suspected
\item Electrocardiogram to rule out cardiac ischemia
\item Pulse oximetry to exclude hypoxia
\end{itemize}

\section*{Treatment Options}
\begin{itemize}
\item Treat underlying cause (anxiety, pulmonary embolus, salicylate overdose)
\item Reassurance and verbal coaching to slow breathing
\item Benzodiazepines for acute panic-related hyperventilation
\item Rebreathing into a mask or reservoir may be used in emergency
\item Oxygen for hypoxic people
\item IV calcium gluconate for severe tetany
\item Cognitive behavioral therapy for chronic anxiety
\end{itemize}

\section*{References}

\begin{thebibliography}{99}
\bibitem{harrison-hypervent} Longo DL, et al. \emph{Harrison's Principles of Internal Medicine}. 21st ed. McGraw-Hill; 2022. Chapter 51: Hyperventilation.
\bibitem{uptodate-hypervent} Schwartzstein RM, et al. Hyperventilation syndrome. In: UpToDate, Post TW (Ed), UpToDate, Waltham, MA; 2025.
\bibitem{trend} Trend P, et al. Hyperventilation syndrome: diagnosis and management. \emph{Clin Med (Lond)}. 2023;23(2):140-145.
\bibitem{bass} Bass C, et al. Hyperventilation syndrome: a clinical review. \emph{J Psychosom Res}. 2022;157:110928.
\bibitem{west} West JB, et al. \emph{West's Respiratory Physiology: The Essentials}. 12th ed. Wolters Kluwer; 2023. Chapter 8: Acid-Base Balance.
\bibitem{balon} Balon R, et al. Panic disorder and hyperventilation. \emph{CNS Spectr}. 2023;28(1):38-45.
\end{thebibliography}
